\documentclass[12pt]{article}

% Idioma e codificação
\usepackage[brazil]{babel}
\usepackage[utf8]{inputenc}
\usepackage[T1]{fontenc}

% Matemática e listas
\usepackage{amsmath, amssymb}
\usepackage{enumitem}

% Layout e cores
\usepackage{geometry}
\usepackage{graphicx}
\usepackage{xcolor}

% Cores institucionais
\definecolor{maincolor}{RGB}{0,100,50}
\definecolor{forestgreen}{RGB}{0,100,50}

% Cabeçalho e rodapé
\usepackage{fancyhdr}

% Blocos
\usepackage[most]{tcolorbox}
\newtcolorbox{block}[1]{%
  colback=green!5,
  colframe=maincolor,
  coltitle=white,
  title=#1,
  fonttitle=\bfseries,
  boxrule=1.5pt,
  arc=3mm,
  left=6pt,
  right=6pt,
  top=6pt,
  bottom=6pt
}

% Margens
\geometry{margin=2.5cm}

% Fancyhdr
\pagestyle{fancy}
\fancyhf{}

\renewcommand{\headrulewidth}{2pt}
\renewcommand{\footrulewidth}{2pt}

\renewcommand{\headrule}{\hbox to\headwidth{%
  \color{forestgreen}\leaders\hrule height \headrulewidth\hfill}}

\renewcommand{\footrule}{\hbox to\headwidth{%
  \color{forestgreen}\leaders\hrule height \footrulewidth\hfill}}

% Altura do cabeçalho (corrigida)
\setlength{\headheight}{52pt}
\addtolength{\topmargin}{-20pt}

% Cabeçalho
\fancyhead[L]{%
  \includegraphics[height=1.2cm]{Universidade_de_Rio_Verde.jpg}
}

\fancyhead[R]{%
  \raggedleft
  \textcolor{forestgreen}{%
    \small
    \textbf{Lista de Exercícios}\\
    Me.~Sergio Souza Novak\\
    \textit{Banco de Dados}
    }
    }
    
    % Rodapé
    \fancyfoot[C]{\textcolor{forestgreen}{\thepage}}
    
    % =========================
    \begin{document}
    % =========================
    
    \vspace*{0.2cm}
    \begin{center}
      {\Large \textbf{Atividade Prática}}\\[0.2cm]
      {\large Banco de dados}\\[0.2cm]
      Tema: Modelo Entidade Relacionamento (MER)
    \end{center}
    
    \vspace{0.5cm}
    
\section*{Atividade: Modelagem Entidade--Relacionamento}

\subsection*{Contexto}

Uma universidade deseja desenvolver um \textbf{sistema de controle acadêmico} para gerenciar informações sobre alunos, cursos, disciplinas, professores e matrículas. O sistema deverá armazenar e organizar os dados de forma consistente, permitindo consultas e relatórios acadêmicos.

\subsection*{Descrição do Sistema}

Considere as seguintes regras de negócio:

\begin{itemize}
  \item Cada \textbf{aluno} possui um número de matrícula único, nome, CPF, data de nascimento e e-mail.
  \item Um \textbf{curso} possui um código identificador, nome, carga horária total e modalidade (presencial ou a distância).
  \item Cada aluno está vinculado a \textbf{um único curso}, mas um curso pode possuir \textbf{vários alunos}.
  \item As \textbf{disciplinas} possuem código, nome e carga horária.
  \item Um curso é composto por \textbf{várias disciplinas}, e uma mesma disciplina pode fazer parte de \textbf{mais de um curso}.
  \item Os \textbf{professores} possuem um identificador, nome, titulação e área de atuação.
  \item Cada disciplina é ministrada por \textbf{um professor}, enquanto um professor pode ministrar \textbf{várias disciplinas}.
  \item O aluno pode se matricular em várias disciplinas, e cada disciplina pode ter vários alunos matriculados.
  \item Para cada matrícula em disciplina, devem ser armazenadas a \textbf{nota final} e a \textbf{frequência} do aluno.
\end{itemize}

\subsection*{Atividades}

Com base na descrição apresentada, responda:

\begin{enumerate}
  \item Identifique as \textbf{entidades} do sistema.
  \item Identifique os \textbf{atributos} de cada entidade, destacando as chaves primárias.
  \item Defina os \textbf{relacionamentos} entre as entidades.
  \item Indique as \textbf{cardinalidades} e participações (mínima e máxima) de cada relacionamento.
  \item Elabore o \textbf{Diagrama Entidade--Relacionamento (DER)}, utilizando a notação vista em aula.
\end{enumerate}

\subsection*{Orientações}

\begin{itemize}
  \item O diagrama pode ser feito à mão (de forma legível) ou utilizando uma ferramenta de modelagem.
  \item Utilize corretamente os conceitos de entidade forte, entidade associativa (se necessário), atributos simples e compostos.
  \item Justifique brevemente as decisões de modelagem, quando julgar necessário.
\end{itemize}

    
    


% =========================
\end{document}
% =========================
